\section{Preserving charts and assembling the embedding}
\label{pad:assembly}

The local constructions now provide a uniform chance of separation at a
pair's distance scale. Two estimates let us use that chance in a complete
schedule: coarser scales rarely destroy the common charts required by the
test, and finer scales move too little label mass to erase its certificate.
We first collect these estimates, then choose the scale spacing and write
the resulting finite cut embedding.

\subsection{The one-step estimates}
All constants in this section depend only on the fixed choices in
Section~\ref{sets:parameters}. In particular they are independent of the
number of columns, band components, chart identifiers, and scales.

\begin{proposition}[Movement and success]\label{pad:success}
There are constants $M<\infty$ and $b,p>0$ such that the following hold
at every step of scale $r$.
\begin{enumerate}[label=(\roman*)]
\item At every point of $X$, in either system, the symmetric-difference
movement of the unrefined set is at most $Mr$.
\item Given any past for which a vertex pair at distance in $[r,2r)$ is
eligible, the step has probability at least $p$ of producing one of these
certificates:
\begin{enumerate}[label=(\alph*)]
\item unrefined symmetric-difference separation at least $br$ in one
system;
\item a non-neighbor tip $e$ and a set of labels of measure at least $br$
in system $A$ on which both pair vertices are active and $e$ is inactive,
with the vertices strictly on opposite sides of that link.
\end{enumerate}
\end{enumerate}
\end{proposition}
\begin{proof}
Lemma~\ref{main:validity} bounds movement in the main clipping.
Secondary movement is bounded by Lemma~\ref{sec:validity}; increments are bounded by
Lemma~\ref{sec:increment-bound}. These estimates apply even at points
outside the finite vertex list. There are no simultaneous main updates
at a point except at cell overlaps, where every update is unchanged.
Secondary updates take a maximum, whose movement is bounded without
summing over groups. Taking the largest of the bounds gives $M$.

Work in scale units. If one of the cutoffs in
\eqref{sets:cutoffs} fails, no change preserves separation at least the
minimum of $10^{-4}a_2,\kappa_AE_A,\kappa_B10^{-4}R$.
Otherwise the interior requirements hold with probability at least $1/4$.
Conditional on the bands and cells, a missing descending path is handled
by Corollary~\ref{main:pair-separation}: keep that endpoint's cell and suppress the other,
giving positive separation with probability at least $1/4$ conditional
on the main mode. If both paths exist, an excluded pair has the component
certificate on a label set of measure at least $\kappa_Aj/4$, which the
no-change mode preserves.

For a surviving pair, admission at an endpoint is handled by
Lemma~\ref{sec:admitted-success}, including its no-change alternative.
If neither orientation admits either endpoint, apply
Lemma~\ref{sec:remaining}. Each case uses one of four modes, at most one
fair orientation, and at most $2M_{\mathrm{pack}}+1$ prescribed group coins. It follows,
for example, that $p=2^{-(2M_{\mathrm{pack}}+10)}$ is a permissible lower bound after
all interior and mode probabilities. Choose $b>0$ below all the finitely
many displayed gains and certificate masses. These choices depend only
on the fixed parameters.
\end{proof}

\subsection{The probability of splitting a small ball}
The following estimate concerns chart identifiers, not cut components.
System $A$ is refined only after all the updates.

\begin{proposition}[Chart persistence]\label{pad:persistence}
There exist constants $C_0<\infty$ and $\rho_0>0$ such that the following
holds. Condition on the past at a step of scale $r$. Suppose a fixed
closed ball $B_d(v,\rho r)$ lies in a single old region in each system.
For $0<\rho\le\rho_0$, the probability that this ball is split between
different new regions in either system is at most
\begin{equation}\label{pad:loss}
 C_0(\rho+\eta_0).
\end{equation}
Here $\eta_0$ is the relative mesh allowance from
Section~\ref{sets:setup}.
\end{proposition}
\begin{proof}
Normalize $r=1$ and write $B=B_d(v,\rho)$. Every point of $B$ has
a shortest path to $v$ contained in $B$. Joining two such paths gives
a path of length at most $2\rho$ between any two points of $B$.
We will use these paths for all band-distance comparisons.

There are two ways an update could give different chart identifiers on
$B$: its domain could end inside the ball, or different formulas could
win there. Strict margins at domain boundaries prevent the first kind
of change. The comparison estimate of Corollary~\ref{int:comparisons}
controls the second.

\paragraph{Height bands.}
The probability that an endpoint of either random height tiling lies
within $\rho$ of $h(v)$ is $O(\rho+\eta_0)$. Exclude this event and
condition on the tilings. The ball lies wholly in one component of a
broad band and one component of a narrow band: all its center paths
are paths in those bands.
For clipping, each such band is either processed in its entirety or left unchanged.
A narrow band not contained in a broad band has no secondary seeds,
so secondary clipping also leaves it unchanged.

\paragraph{Main clipping.}
Condition further on the radial cells and their coins. Suppose first
that $B$ meets more than one width-$D$ cell, including through different
representatives of a quotient point. For any touched cell $C$ and any
$w\in B\cap C$, a path through the center to a point in another cell
reaches a frontier point $z\in B\cap\partial C$ within length $2\rho$.
If $C$ is available, its old comparison field is the common old field
on $B$. Thus
\[
 F_C(w)\le s(z)+(1-\zeta_A)d_{J_1}(w,z)-3E_A
          \le s(w)+4\rho-3E_A.
\]
For $\rho_0<E_A/10$ the proposal cannot raise the field anywhere on
$B\cap C$. Unavailable cells are unchanged as well. Hence the whole
ball retains its old identifier. In particular we do not need to
exclude a random radial boundary crossing.

It remains to consider a ball lying in one cell $C$. If that cell is
unavailable or its coin is suppressed, it is unchanged. Otherwise its
comparison expression on $B$ is
\[
 \max\{s,\min\{F_C,p_{h,C}\}\}.
\]
Here $F_C=\widetilde p_C-4E_A$ is one interpolation construction,
$p_{h,C}$ is a fixed height function plus its independent offset, and
$s$ is the only branch with no fresh noise. The input $p_C$, portals,
and bands have all been fixed before these noises are sampled.
Every center path has length at most $\rho$ in the interpolation band
and in the old chart band. Corollary~\ref{int:comparisons}, with the
fixed Lipschitz bound of the height function, therefore gives one
winning branch and, if that branch is the interpolation, one named atom throughout $B$,
except with conditional probability $O(\rho+\eta_0)$.
A common old branch retains its identifier; a common spatial atom or
height branch gives a common new chart. At interpolation layer seams
the inherited atom keeps its name, as required by the interpolation
theorem.

\paragraph{Secondary clipping.}
Condition on the cells, surviving groups, comparison orientations,
admissions, and group coins. These choices precede the fresh clipping
noises. A group whose domain meets $B$ has an admitted seed at band
distance at most $R+2\rho<4R$ from $v$. Proposition~\ref{pack:bound}
therefore bounds the number of such groups by $M_{\mathrm{pack}}$.

If a group's domain meets but does not contain $B$, its boundary is
reached within length $2\rho$ from every point of its domain in the
ball. Lemma~\ref{sec:validity} puts the proposal below the old field
at this spatial boundary by a fixed positive margin; height-boundary
crossings have already been excluded. The fixed Lipschitz constants
of the proposal branches and the old field extend this nonaction to
the domain's intersection with the ball when $\rho_0$ is small enough.
Outside the domain there is no update by definition. Discard such a group. Every
remaining relevant domain contains $B$, so its proposal is described
by one fixed expression throughout the ball.

Each relevant domain lies in its seed's old chart and meets $B$.
Since regions partition $X$, that chart must be the ball's common old
chart. Thus all these proposals use the same old comparison orientation
and label variables. The ball lies in one narrow-band component, so
each group contributes only the interpolation for that component. Their
combined expression is the maximum of $p_*$ and the group proposals,
each proposal being the minimum of $T_F$, $H_F$, and, in an ordinary
chart, $O_F$. The $T_F$ are signed, shifted interpolation functions;
the $H_F$ and $O_F$ have independent additive offsets. The only input
without fresh noise is the old field $p_*$. The group data and the
interpolation inputs are already fixed, and fresh noise families are
independent between groups. All center paths are permitted in their
narrow-band metrics and in the old chart band.

Apply Corollary~\ref{int:comparisons} to this expression. Each
interpolation contributes its bounded list of candidate atom names;
the number of lists is bounded by $M_{\mathrm{pack}}$. All constants
are therefore uniform in the graph and the past. Except with
conditional probability $O(\rho+\eta_0)$, one branch and, where
applicable, one atom name wins throughout $B$. The resulting chart
identifier is common to the ball. Repeated appearances of an inherited
atom use its interpolation and original atom name, so crossing a
portal collar does not create another identifier.

No-change and increment modes introduce no chart identifiers. Taking
$\rho_0$ smaller than the finitely many bounds used above and summing
the exceptional probabilities proves~\eqref{pad:loss}.
\end{proof}

\subsection{Choice of spacing and the final map}
Choose $Q=2^m$ large enough for the support condition in
Lemma~\ref{sec:increment-bound}, for $20/Q<\rho_0$, and for
\begin{equation}\label{pad:spacing}
 \frac{40C_0}{Q-1}<\frac12,
 \qquad \frac{6M}{Q-1}<\frac b2.
\end{equation}
Choose the mesh allowance so that $\eta_0\le r_f/r$ for every
pair of scales $r_f<r$ in the finite collection of schedules. Then the
mesh term in \eqref{pad:loss} is no larger than the corresponding scale
ratio. This choice does not change any of the constants above.

Fix a pair and its designated scale $r_f$, and consider
$B_d(x,20r_f)$. Initially it has one chart in each system. Conditional
on having retained that property up to a coarser scale $r$, its first
loss probability is at most $40C_0 r_f/r$. Therefore
\[
 \Pr[\text{a first loss before }r_f]
 \le40C_0\sum_{k\ge1}Q^{-k}<\frac12.
\]
On the complementary event the pair is eligible at scale $r_f$.
Furthermore any $D$-cell needed for its main update has its entire
$1$-neighborhood inside this ball: its weak diameter is at most $10D$.
The old chart bands contain the short paths used in all tests.

By Proposition~\ref{pad:success}, conditional on this event the scale
produces a certificate with probability at least $p$. Subsequent
unrefined movement at any one point is deterministically at most
\[
 M\sum_{k\ge1}r_fQ^{-k}=\frac{Mr_f}{Q-1}.
\]
An ordinary pair separation loses at most twice this quantity. For a
component certificate, remove labels whose membership later changes
at either vertex or at its tip; their total measure is at most three
times this quantity. At least $br_f/2$ of the original certificate
therefore remains. Lemma~\ref{sets:component-refinement} gives final
separation for all those labels, regardless of changes elsewhere.
In either case the final sum of the refined $A$ distance and the
unrefined $B$ distance is at least $br_f/2$ on the successful event.
Its expectation in the designated schedule is consequently at least
\[
 \frac12\,p\,\frac{br_f}{2}
 \ge\frac{pb}{8}d(x,y).
\]

The designated schedule is the residue class containing the dyadic
exponent of $r_f$; every other schedule contributes a nonnegative
distance. Mix uniformly over the $m$ schedules. Within each schedule use the
probability law of the random construction, weighting each outcome by
its probability. More precisely, an outcome is a branch of the finite
execution tree, and its weight is the product of the conditional
probabilities of the choices along that branch. These weights sum to one
even though later geometric choices depend on earlier outcomes.
Take the finite direct sum of their measure spaces
with these weights, and of the two systems. Subtract the
assignment at a fixed reference vertex in every summand. By
Lemma~\ref{sets:component-refinement} the resulting map $F$ is into a
real $L_1$ space and satisfies
\[
 \frac{pb}{8m}d(x,y)\le\|F(x)-F(y)\|_1\le2d(x,y)
 \qquad(x,y\in V).
\]
This map can be written using finitely many cut coordinates. For every
final outcome and every label in system $B$, record the subset of $V$
that is active. For system $A$, record separately, for each active column
component, the subset of original vertices active in that component.
Each such subset $S$ contributes the cut distance
$|\1_S(x)-\1_S(y)|$. Integrating over the labels and averaging with the
schedule and outcome weights therefore gives numbers $w_S\ge0$ such that
\[
 \|F(x)-F(y)\|_1
   =\sum_{\varnothing\ne S\subsetneq V}
       w_S|\1_S(x)-\1_S(y)|.
\]
Empty and full subsets contribute zero and have been discarded. Every
remaining weight is finite: if $x\in S$ and $y\notin S$, its contribution
is bounded by the already established upper bound $2d(x,y)$. Hence
\[
 v\longmapsto\bigl(w_S\1_S(v)\bigr)_{
                         \varnothing\ne S\subsetneq V}
\]
is a map into a finite-dimensional real $\ell_1$ space with exactly the
same distances as $F$.

The column realization identifies these distances with the original
$d_G$. Rescale $F$ by $8m/(pb)$ to obtain Theorem~\ref{thm:main} with
$C=16m/(pb)$. If $V$ has a single point use the zero map. Every constant
is independent of the graph size and of its positive edge lengths.
